\documentclass[11pt]{article}
\usepackage[a4paper,margin=26mm]{geometry}
\usepackage{amsmath,amssymb,mathtools}
\usepackage{unicode-math}
\setmainfont{Latin Modern Roman}
\setmathfont{Latin Modern Math}
\usepackage{xurl,hyperref}
\hypersetup{hidelinks,pdftitle={UnsolvedMath Problem 10400030}}
\newcommand{\sref}[2]{\hyperlink{source:#1:#2}{[S#2]}}
\newcommand{\cataloguescope}[1]{#1}
\begin{document}
% UnsolvedMath ID 10400030; source catalogue code AMR-103-0030
\setcounter{section}{10400029}
\section{Question 2.8 — (T.}\label{10400030}
{\small\sffamily UnsolvedMath ID 10400030 · Topology}
\subsection{Definitions and mathematical statement}

\paragraph{Shared notation.}$\mathbb N=\{1,2,\ldots\}$, and $\mathbb Z,\mathbb Q,\mathbb R,\mathbb C$ denote the integers, rationals, reals and complex numbers. Any other conventions are specified locally.


A knot or link is a smooth oriented embedding of one or more circles in the oriented $3$-sphere, considered up to orientation-preserving ambient isotopy. Write $O$ for the unknot, which bounds an embedded disk. Let $R\mathcal K$ be the free $R$-module on oriented knot isotopy classes. Extend a knot invariant to singular knots with ordinary transverse double points by $v(K_\times)=v(K_+)-v(K_-)$. Let $F_d(R\mathcal K)$ be the submodule generated by the alternating resolution sums of knots with $d$ such points. An $R$-valued finite-type (Vassiliev) invariant of degree at most $d$ is an $R$-linear map $R\mathcal K\to R$ vanishing on $F_{d+1}(R\mathcal K)$. For a finite group $G$, put $h_G(K)=\#\operatorname{Hom}(\pi_1(S^3\setminus K),G)$, counting homomorphisms themselves rather than conjugacy classes. A finite quandle is a finite set $X$ with a binary operation $\triangleright$ that is idempotent, has bijective right translations, and satisfies right self-distributivity. The knot quandle is the quandle of oriented meridians with its usual conjugation operation, or equivalently the quandle presented by arc labels and crossing relations. Put $h_X(K)=\#\operatorname{Hom}_{\mathrm{quandle}}(Q(K),X)$.
\paragraph{Statement recorded in the supplied source.}
For finite groups $G$ other than dihedral groups, can $h_G$ be approximated pointwise by finite-type invariants? Seek one sequence $(v_j)$ of real-valued finite-type invariants, with degrees allowed to grow, for which $v_j(K)\to h_G(K)$ for every knot $K$. Investigate also the source's weaker query whether finite-type invariants convey any information about $h_G$ if full approximation fails, and its refinement replacing $h_G$ by $h_X$ for finite quandles $X$. \sref{10400030}{2}
\subsection{Short English statement}
Can finite-type invariants approximate finite-group representation counts of knot groups, and the analogous finite-quandle coloring counts; if not, can they still provide information about those counts?
\subsection{Sources}
\begin{description}
\item[\hypertarget{source:10400030:1}{S1}] ULAM. AI and the UnsolvedMath Contributors, \emph{UnsolvedMath: A Curated Collection of Open Mathematics Problems}, record 10400030 (AMR-103-0030). CC BY 4.0. \url{https://huggingface.co/datasets/ulamai/UnsolvedMath}.
\item[\hypertarget{source:10400030:2}{S2}] T. Ohtsuki (editor), \emph{Problems on invariants of knots and 3-manifolds}, Geometry \& Topology Monographs 4 (2002), pp. 377--572 (published 2004), item 2.8 and its surrounding definitions and remarks. \url{https://arxiv.org/abs/math/0406190v1}.
\end{description}
\label{end:10400030}
\end{document}
